\section{引言}
\label{sec:intro}

卖方深度研报是二级市场投研最主要的文字产品：分析师读完年报、季报与公告，搜集行业与竞争资料，建立盈利预测和估值模型，最后写成二三十页、数十张图表的报告。一份年报点评加估值模型通常要花费分析师数天时间。大语言模型在长文档理解与写作上的能力，使“让模型直接读年报写研报”看起来触手可及。

然而在我们的早期实验中，模型直接生成的研报存在三类系统性失败，每一类都足以让整篇报告失去可信度：

\begin{itemize}
  \item \textbf{F1 数字幻觉。}模型会记错或编造财务数字，混淆合并报表与母公司报表、本期与上期、营业收入与营业总收入。即使数字抄对，也看不出它来自年报哪一页。
  \item \textbf{F2 计算不可靠。}同比增速、占比、单季拆分、三表勾稽和折现估值由模型“心算”时会出错，且结果随采样波动，无法复现，也无法做敏感性分析。
  \item \textbf{F3 外部信息无出处。}年报之外的行业规模、市场份额、一致预期、竞争对手动态，模型倾向于凭参数记忆补全，既可能过时，也可能是虚构。
\end{itemize}

第四个困难来自评价本身（\textbf{F4}）：研报的好坏涉及分析是否成立、论证是否充分、表达是否专业，长期依赖资深分析师的主观判断。没有可复核的评价，就无法判断系统改动是否真的带来改进，也无法驱动自动化的迭代。

\paragraph{本文的立场。}我们不试图让模型“更会算”，而是把计算从模型手里拿走。\sysname{} 的总原则可以概括为一句话：\textbf{模型读、抄、写、判断；程序校验、计算、拦截。}模型给出的每一个数字，要么能在原文里逐字找到，要么由程序从真值表算出，要么登记在外部来源表中并在句中标注来源编号；否则会被程序拒收或拦下。在此基础上，研报的生产被拆成若干只靠文件交接、互不继承上下文的会话，每一步都有程序闸门；研报的评价被拆成程序判定的数据闭合与评审模型判定的分析质量两部分，评审模型本身也要先经过校准。

\paragraph{贡献。}本文的贡献如下：
\begin{enumerate}
  \item \textbf{理解与计算分离的系统架构}（第 \ref{sec:principles}、\ref{sec:arch} 节）。我们把九条设计规范写成可由程序检查的不变量，并给出每条规范的执行机制与测试，使“每个数字可追溯”从口号变为流水线的硬约束。
  \item \textbf{面向真实年报版式的证据约束抽取}（第 \ref{sec:extraction} 节）。针对 A 股年报中的稀疏网格、跨页续表、科目折行、同页多单位声明等版式问题设计规则通道，并与模型通道交叉校验；两通道冲突时记为问题交人工，不静默取舍。
  \item \textbf{数字护栏与数字锁}（第 \ref{sec:guards} 节）。护栏按每个数字的展示精度设定容差，只在结论性产物中找出处，从而压低“碰巧放行”的概率；数字锁保证改写只能删减、不能新增或改动数字；二者共同把 LLM 改写报告的风险限制在措辞与结构层面。
  \item \textbf{先搜集、后分析的多会话研报流水线}（第 \ref{sec:pipeline} 节）。五个角色（研究助理、分析师、主笔、复核编辑、排版设计师）在仓库之外的隔离工作区中顺序工作，11 道程序闸门逐阶段核对产物及其新鲜度。
  \item \textbf{可校准的评测闭环}（第 \ref{sec:eval} 节）。四层评分标准（硬闸门、数据闭合、分析质量、表达与合规）加负向条目；评分标准对生成方不可见，问题单按类型回传；我们用重复评审、程序制造缺陷和人工复核三组证据检验评审模型是否可信。
\end{enumerate}

\paragraph{结果概览。}在贵州茅台、宁德时代、招商轮船三家 2025 年年报上，规则通道覆盖 314 个人工核验值中的 293 个（改进前为 128 个），错配为 0（第 \ref{sec:exp-extract} 节）。流水线在“先搜集外部资料”的基础上加入深度专题调研后，中际旭创与宇树科技两份研报的质检得分由 74、67 分提升至 90、88 分，均达到交付标准（第 \ref{sec:exp-quality} 节）。评审模型的校准显示：重复盲评逐条一致率 86\%～91\%、总分极差 7～8 分；无来源数字、机器腔、评级矛盾等缺陷能够判出，整节删除类缺陷的判出率较低，原因分析见第 \ref{sec:exp-calib} 节。我们同时如实报告这些数字的局限：样本量小、版本之间不是对照实验、评审与生成模型同族（第 \ref{sec:limits} 节）。
